\section{Introdução}
\label{sec:intro}

\subsection{Elementos de interface}

Em muitos problemas de engenharia o comportamento que mais importa acontece em superfícies de descontinuidade:
fraturas e falhas geológicas, juntas em maciços rochosos, contatos entre materiais distintos, caminhos de propagação de
trincas. No método dos elementos finitos, uma forma clássica de representar essas descontinuidades é inserir, ao longo
delas, \emph{elementos de interface} --- também chamados de elementos de espessura nula ou, quando associados a uma lei
de tração--separação, elementos coesivos~\cite{paulino}.

Em 2D, um elemento de interface linear é um quadrilátero de quatro nós cujos nós vêm aos pares geometricamente
coincidentes (\cref{fig:ie-conceito}). Dois nós pertencem a um lado da descontinuidade e dois ao outro. Como a espessura
é nula, a grandeza cinemática relevante deixa de ser a deformação e passa a ser o \emph{salto de deslocamento} entre os
dois lados, $[\![\mathbf{u}]\!] = \mathbf{u}^{+} - \mathbf{u}^{-}$, que mede abertura e deslizamento da interface.

\begin{figure}[ht]
\centering
\begin{tikzpicture}[scale=1.05, every node/.style={font=\small}]
  % (a) contínuo
  \begin{scope}
    \filldraw[cefill] (0,0) rectangle (3,1); \filldraw[cefill] (0,1) rectangle (3,2);
    \draw[seledge] (0,1) -- (3,1);
    \node[elemlabel] at (1.5,1.5) {U}; \node[elemlabel] at (1.5,0.5) {L};
    \node[vreg, minimum size=0.38cm] at (0,1) {1}; \node[vreg, minimum size=0.38cm] at (3,1) {2};
    \node at (1.5,-0.5) {(a) malha contínua};
  \end{scope}
  % (b) com elemento de interface (espessura exagerada)
  \begin{scope}[xshift=4.2cm]
    \filldraw[cefill] (0,0) rectangle (3,1); \filldraw[iefill] (0,1) rectangle (3,1.35);
    \filldraw[cefill] (0,1.35) rectangle (3,2.35);
    \node[elemlabel] at (1.5,1.85) {U}; \node[elemlabel] at (1.5,0.5) {L};
    \node[font=\scriptsize\bfseries, text=orange!40!black] at (1.5,1.175) {I};
    \node[vie, minimum size=0.36cm] at (0,1) {1}; \node[vie, minimum size=0.36cm] at (3,1) {2};
    \node[vie, minimum size=0.36cm] at (3,1.35) {3}; \node[vie, minimum size=0.36cm] at (0,1.35) {4};
    \draw[decorate, decoration={brace, amplitude=3pt}] (3.3,1.4) -- (3.3,0.95)
      node[midway, right=3pt, align=left, font=\scriptsize] {espessura nula\\(exagerada)};
    \node at (1.5,-0.5) {(b) elemento de interface};
  \end{scope}
  % (c) abertura
  \begin{scope}[xshift=9.6cm]
    \filldraw[cefill] (0,0) rectangle (3,1);
    \filldraw[iefill] (0,1) -- (3,1) -- (3.25,1.55) -- (0.15,1.3) -- cycle;
    \filldraw[cefill] (0.15,1.3) -- (3.25,1.55) -- (3.25,2.55) -- (0.15,2.3) -- cycle;
    \node[elemlabel] at (1.7,1.95) {U}; \node[elemlabel] at (1.5,0.5) {L};
    \draw[-{Stealth[length=4pt]}, red!70!black, line width=0.8pt] (3,1) -- (3.25,1.55);
    \draw[-{Stealth[length=4pt]}, red!70!black, line width=0.8pt] (0,1) -- (0.15,1.3);
    \node[text=red!70!black, font=\scriptsize] at (4.05,1.25) {$[\![\mathbf{u}]\!]=\mathbf{u}^{+}-\mathbf{u}^{-}$};
    \node at (1.6,-0.5) {(c) abertura da interface};
  \end{scope}
\end{tikzpicture}
\caption{O elemento de interface. (a) Dois elementos contínuos compartilham uma aresta (em vermelho). (b) A aresta é
\emph{aberta}: os nós 1 e 2 são duplicados em 4 e 3, e entre U e L surge um quadrilátero de quatro nós com
$\mathbf{x}_1=\mathbf{x}_4$ e $\mathbf{x}_2=\mathbf{x}_3$ (desenhado com espessura exagerada). (c) Durante a simulação os
dois lados podem se separar; o elemento mede o salto de deslocamento.}
\label{fig:ie-conceito}
\end{figure}

Há duas estratégias para colocar esses elementos na malha. Na estratégia \emph{a priori}, todos os elementos de
interface são criados antes da análise, ao longo de todas as superfícies em que uma descontinuidade pode ocorrer. Na
estratégia \emph{adaptativa}, eles são inseridos durante a simulação, apenas onde e quando um critério (de tensão, de
dano etc.) indica que a descontinuidade se formou~\cite{paulino}. A segunda estratégia evita a rigidez artificial e o
custo de milhares de elementos inúteis, mas exige um algoritmo que altere a topologia da malha de forma incremental,
local e eficiente, sobre uma malha que \emph{já contém} elementos de interface.

\subsection{O problema da inserção}
\label{sec:problema}

\begin{nota}[Problema]
\textbf{Dados:} uma malha 2D orientada $\mathcal{M}$ e um conjunto $S$ de arestas internas (cada uma compartilhada por
dois elementos).\\
\textbf{Obter:} uma malha $\mathcal{M}'$ em que cada aresta de $S$ foi \emph{aberta}: os dois elementos que a
compartilhavam passam a se ligar através de um novo elemento de interface de quatro nós, e os vértices são duplicados
apenas o necessário para que a topologia continue consistente.
\end{nota}

A \cref{fig:exemplo-intro} mostra a entrada e a saída para o exemplo usado ao longo de todo o texto: uma malha híbrida
de quadriláteros e triângulos, com 12 arestas selecionadas e numeradas na ordem de inserção. Parece simples, mas há
várias dificuldades:
\begin{enumerate}
\item \textbf{Quantas cópias de cada vértice?} Um vértice interno tocado por uma única aresta selecionada não precisa
ser duplicado geometricamente (a trinca termina nele). Com duas arestas, ele vira dois. Com três ou mais, surge uma
\emph{junção} (em T, em X, em estrela), e os elementos de interface não conseguem fechar a vizinhança sozinhos:
aparece um \emph{buraco} poligonal, visível como as regiões hachuradas na \cref{fig:exemplo-intro}b.
\item \textbf{Quais elementos vão para cada cópia?} A estrela do vértice precisa ser dividida em setores delimitados
pelas arestas selecionadas, e cada setor deve passar a usar uma cópia diferente.
\item \textbf{Borda.} Na borda da malha a abertura não precisa ser fechada: o vértice sempre ganha uma cópia a mais.
\item \textbf{Incrementalidade.} Em uma análise adaptativa as arestas chegam aos poucos e em qualquer ordem. O
resultado não deveria depender dessa ordem.
\item \textbf{Estruturas compactas.} Estruturas baseadas em índices, como a CHE e a IBHM~\cite{che,ibhm}, armazenam a
malha em vetores de inteiros, sem ponteiros. Os próprios autores da IBHM apontam como limitação o uso em remalhamento
dinâmico~\cite{ibhm}. Inserir elementos é exatamente esse tipo de alteração.
\item \textbf{Extensão para 3D.} O objetivo final é inserir elementos em faces de malhas 3D. As decisões tomadas em 2D
devem ser transportáveis.
\end{enumerate}

\begin{figure}[ht]
\centering
\begin{subfigure}[b]{0.42\linewidth}\centering
  \begin{tikzpicture}[scale=1.250, every node/.style={font=\scriptsize}]
\filldraw[cefill] (0.000,0.000) -- (1.000,0.000) -- (1.000,1.000) -- (0.000,1.000) -- cycle;
\filldraw[cefill] (1.000,0.000) -- (2.000,0.000) -- (2.000,1.000) -- (1.000,1.000) -- cycle;
\filldraw[cefill] (2.000,0.000) -- (3.000,0.000) -- (3.000,1.000) -- (2.000,1.000) -- cycle;
\filldraw[cefill] (0.000,1.000) -- (1.000,1.000) -- (1.000,2.000) -- (0.000,2.000) -- cycle;
\filldraw[cefill] (1.000,1.000) -- (2.000,1.000) -- (2.000,2.000) -- (1.000,2.000) -- cycle;
\filldraw[cefill] (2.000,1.000) -- (3.000,1.000) -- (3.000,2.000) -- (2.000,2.000) -- cycle;
\filldraw[cefill] (3.000,1.000) -- (4.000,1.000) -- (3.000,2.000) -- cycle;
\filldraw[cefill] (4.000,1.000) -- (4.000,2.000) -- (3.000,2.000) -- cycle;
\filldraw[cefill] (0.000,2.000) -- (1.000,2.000) -- (1.000,3.000) -- (0.000,3.000) -- cycle;
\filldraw[cefill] (1.000,2.000) -- (2.000,2.000) -- (2.000,3.000) -- (1.000,3.000) -- cycle;
\filldraw[cefill] (2.000,2.000) -- (3.000,2.000) -- (3.000,3.000) -- (2.000,3.000) -- cycle;
\filldraw[cefill] (3.000,2.000) -- (4.000,2.000) -- (4.000,3.000) -- cycle;
\filldraw[cefill] (3.000,2.000) -- (4.000,3.000) -- (3.000,3.000) -- cycle;
\draw[seledge] (1.000,2.000) -- (1.000,3.000);
\node[edgenum] at (0.840,2.500) {1};
\draw[seledge] (1.000,2.000) -- (2.000,2.000);
\node[edgenum] at (1.500,2.160) {2};
\draw[seledge] (2.000,1.000) -- (2.000,0.000);
\node[edgenum] at (2.160,0.500) {3};
\draw[seledge] (3.000,2.000) -- (3.000,1.000);
\node[edgenum] at (3.160,1.500) {4};
\draw[seledge] (3.000,3.000) -- (3.000,2.000);
\node[edgenum] at (3.160,2.500) {5};
\draw[seledge] (2.000,2.000) -- (3.000,2.000);
\node[edgenum] at (2.500,2.160) {6};
\draw[seledge] (2.000,2.000) -- (2.000,1.000);
\node[edgenum] at (2.160,1.500) {7};
\draw[seledge] (1.000,2.000) -- (1.000,1.000);
\node[edgenum] at (1.160,1.500) {8};
\draw[seledge] (0.000,2.000) -- (1.000,2.000);
\node[edgenum] at (0.500,2.160) {9};
\draw[seledge] (3.000,2.000) -- (4.000,2.000);
\node[edgenum] at (3.500,2.160) {10};
\draw[seledge] (3.000,2.000) -- (4.000,1.000);
\node[edgenum] at (3.613,1.613) {11};
\draw[seledge] (3.000,2.000) -- (4.000,3.000);
\node[edgenum] at (3.387,2.613) {12};
\draw[hearrow, black!55] (0.220,0.100) -- (0.780,0.100);
\draw[hearrow, black!55] (0.900,0.220) -- (0.900,0.780);
\draw[hearrow, black!55] (0.780,0.900) -- (0.220,0.900);
\draw[hearrow, black!55] (0.100,0.780) -- (0.100,0.220);
\draw[hearrow, black!55] (1.220,0.100) -- (1.780,0.100);
\draw[hearrow, black!55] (1.900,0.220) -- (1.900,0.780);
\draw[hearrow, black!55] (1.780,0.900) -- (1.220,0.900);
\draw[hearrow, black!55] (1.100,0.780) -- (1.100,0.220);
\draw[hearrow, black!55] (2.220,0.100) -- (2.780,0.100);
\draw[hearrow, black!55] (2.900,0.220) -- (2.900,0.780);
\draw[hearrow, black!55] (2.780,0.900) -- (2.220,0.900);
\draw[hearrow, black!55] (2.100,0.780) -- (2.100,0.220);
\draw[hearrow, black!55] (0.220,1.100) -- (0.780,1.100);
\draw[hearrow, black!55] (0.900,1.220) -- (0.900,1.780);
\draw[hearrow, black!55] (0.780,1.900) -- (0.220,1.900);
\draw[hearrow, black!55] (0.100,1.780) -- (0.100,1.220);
\draw[hearrow, black!55] (1.220,1.100) -- (1.780,1.100);
\draw[hearrow, black!55] (1.900,1.220) -- (1.900,1.780);
\draw[hearrow, black!55] (1.780,1.900) -- (1.220,1.900);
\draw[hearrow, black!55] (1.100,1.780) -- (1.100,1.220);
\draw[hearrow, black!55] (2.220,1.100) -- (2.780,1.100);
\draw[hearrow, black!55] (2.900,1.220) -- (2.900,1.780);
\draw[hearrow, black!55] (2.780,1.900) -- (2.220,1.900);
\draw[hearrow, black!55] (2.100,1.780) -- (2.100,1.220);
\draw[hearrow, black!55] (3.175,1.089) -- (3.735,1.089);
\draw[hearrow, black!55] (3.709,1.149) -- (3.149,1.709);
\draw[hearrow, black!55] (3.089,1.735) -- (3.089,1.175);
\draw[hearrow, black!55] (3.911,1.265) -- (3.911,1.825);
\draw[hearrow, black!55] (3.825,1.911) -- (3.265,1.911);
\draw[hearrow, black!55] (3.291,1.851) -- (3.851,1.291);
\draw[hearrow, black!55] (0.220,2.100) -- (0.780,2.100);
\draw[hearrow, black!55] (0.900,2.220) -- (0.900,2.780);
\draw[hearrow, black!55] (0.780,2.900) -- (0.220,2.900);
\draw[hearrow, black!55] (0.100,2.780) -- (0.100,2.220);
\draw[hearrow, black!55] (1.220,2.100) -- (1.780,2.100);
\draw[hearrow, black!55] (1.900,2.220) -- (1.900,2.780);
\draw[hearrow, black!55] (1.780,2.900) -- (1.220,2.900);
\draw[hearrow, black!55] (1.100,2.780) -- (1.100,2.220);
\draw[hearrow, black!55] (2.220,2.100) -- (2.780,2.100);
\draw[hearrow, black!55] (2.900,2.220) -- (2.900,2.780);
\draw[hearrow, black!55] (2.780,2.900) -- (2.220,2.900);
\draw[hearrow, black!55] (2.100,2.780) -- (2.100,2.220);
\draw[hearrow, black!55] (3.265,2.089) -- (3.825,2.089);
\draw[hearrow, black!55] (3.911,2.175) -- (3.911,2.735);
\draw[hearrow, black!55] (3.851,2.709) -- (3.291,2.149);
\draw[hearrow, black!55] (3.149,2.291) -- (3.709,2.851);
\draw[hearrow, black!55] (3.735,2.911) -- (3.175,2.911);
\draw[hearrow, black!55] (3.089,2.825) -- (3.089,2.265);
\node[vreg, minimum size=0.41cm] at (0.000,0.000) {$A$};
\node[vreg, minimum size=0.41cm] at (1.000,0.000) {$B$};
\node[vreg, minimum size=0.41cm] at (2.000,0.000) {$C$};
\node[vreg, minimum size=0.41cm] at (3.000,0.000) {$D$};
\node[vreg, minimum size=0.41cm] at (0.000,1.000) {$E$};
\node[vreg, minimum size=0.41cm] at (1.000,1.000) {$F$};
\node[vreg, minimum size=0.41cm] at (2.000,1.000) {$G$};
\node[vreg, minimum size=0.41cm] at (3.000,1.000) {$H$};
\node[vreg, minimum size=0.41cm] at (4.000,1.000) {$I$};
\node[vreg, minimum size=0.41cm] at (0.000,2.000) {$J$};
\node[vreg, minimum size=0.41cm] at (1.000,2.000) {$K$};
\node[vreg, minimum size=0.41cm] at (2.000,2.000) {$L$};
\node[vreg, minimum size=0.41cm] at (3.000,2.000) {$M$};
\node[vreg, minimum size=0.41cm] at (4.000,2.000) {$N$};
\node[vreg, minimum size=0.41cm] at (0.000,3.000) {$P$};
\node[vreg, minimum size=0.41cm] at (1.000,3.000) {$Q$};
\node[vreg, minimum size=0.41cm] at (2.000,3.000) {$R$};
\node[vreg, minimum size=0.41cm] at (3.000,3.000) {$S$};
\node[vreg, minimum size=0.41cm] at (4.000,3.000) {$T$};

\end{tikzpicture}

  \caption{Entrada: malha com as arestas selecionadas (em vermelho) e a ordem de inserção.}
\end{subfigure}\hfill
\begin{subfigure}[b]{0.55\linewidth}\centering
  \begin{tikzpicture}[scale=1.900, every node/.style={font=\scriptsize}]
\fill[holefill] (0.880,2.120) -- (1.120,2.120) -- (1.120,1.880) -- (0.880,1.880) -- cycle;
\fill[holefill] (1.880,1.880) -- (2.000,2.170) -- (2.120,1.880) -- cycle;
\fill[holefill] (3.076,1.848) -- (2.880,1.880) -- (2.880,2.120) -- (3.076,2.152) -- (3.152,2.076) -- (3.152,1.924) -- cycle;
\filldraw[cefill] (0.000,0.000) -- (1.000,0.000) -- (1.000,1.000) -- (0.000,1.000) -- cycle;
\filldraw[cefill] (1.000,0.000) -- (1.880,0.120) -- (1.830,1.000) -- (1.000,1.000) -- cycle;
\filldraw[cefill] (2.120,0.120) -- (3.000,0.000) -- (2.830,1.000) -- (2.170,1.000) -- cycle;
\filldraw[cefill] (0.000,1.000) -- (1.000,1.000) -- (0.880,1.880) -- (0.120,1.880) -- cycle;
\filldraw[cefill] (1.000,1.000) -- (1.830,1.000) -- (1.880,1.880) -- (1.120,1.880) -- cycle;
\filldraw[cefill] (2.170,1.000) -- (2.830,1.000) -- (2.880,1.880) -- (2.120,1.880) -- cycle;
\filldraw[cefill] (3.120,1.120) -- (3.848,1.076) -- (3.076,1.848) -- cycle;
\filldraw[cefill] (3.924,1.152) -- (3.880,1.880) -- (3.152,1.924) -- cycle;
\filldraw[cefill] (0.120,2.120) -- (0.880,2.120) -- (0.880,2.880) -- (0.000,3.000) -- cycle;
\filldraw[cefill] (1.120,2.120) -- (2.000,2.170) -- (2.000,3.000) -- (1.120,2.880) -- cycle;
\filldraw[cefill] (2.000,2.170) -- (2.880,2.120) -- (2.880,2.880) -- (2.000,3.000) -- cycle;
\filldraw[cefill] (3.152,2.076) -- (3.880,2.120) -- (3.924,2.848) -- cycle;
\filldraw[cefill] (3.076,2.152) -- (3.848,2.924) -- (3.120,2.880) -- cycle;
\filldraw[iefill] (1.120,2.120) -- (1.120,2.880) -- (0.880,2.880) -- (0.880,2.120) -- cycle;
\filldraw[iefill] (1.120,1.880) -- (1.880,1.880) -- (2.000,2.170) -- (1.120,2.120) -- cycle;
\filldraw[iefill] (1.830,1.000) -- (1.880,0.120) -- (2.120,0.120) -- (2.170,1.000) -- cycle;
\filldraw[iefill] (2.880,1.880) -- (2.830,1.000) -- (3.120,1.120) -- (3.076,1.848) -- cycle;
\filldraw[iefill] (2.880,2.880) -- (2.880,2.120) -- (3.076,2.152) -- (3.120,2.880) -- cycle;
\filldraw[iefill] (2.120,1.880) -- (2.880,1.880) -- (2.880,2.120) -- (2.000,2.170) -- cycle;
\filldraw[iefill] (1.880,1.880) -- (1.830,1.000) -- (2.170,1.000) -- (2.120,1.880) -- cycle;
\filldraw[iefill] (0.880,1.880) -- (0.947,1.046) -- (1.050,1.049) -- (1.120,1.880) -- cycle;
\filldraw[iefill] (0.120,1.880) -- (0.880,1.880) -- (0.880,2.120) -- (0.120,2.120) -- cycle;
\filldraw[iefill] (3.152,1.924) -- (3.880,1.880) -- (3.880,2.120) -- (3.152,2.076) -- cycle;
\filldraw[iefill] (3.076,1.848) -- (3.848,1.076) -- (3.924,1.152) -- (3.152,1.924) -- cycle;
\filldraw[iefill] (3.152,2.076) -- (3.924,2.848) -- (3.848,2.924) -- (3.076,2.152) -- cycle;
\draw[red, line width=1.1pt] (0.947,1.046) -- (1.050,1.049);
\node[vreg, dot] at (0.000,0.000) {};
\node[vreg, dot] at (1.000,0.000) {};
\node[vie, dot] at (2.120,0.120) {};
\node[vreg, dot] at (3.000,0.000) {};
\node[vreg, dot] at (0.000,1.000) {};
\node[vcol, dot] at (1.000,1.000) {};
\node[vie, dot] at (1.830,1.000) {};
\node[vie, dot] at (3.120,1.120) {};
\node[vie, dot] at (3.924,1.152) {};
\node[vie, dot] at (0.120,1.880) {};
\node[vie, dot] at (0.880,2.120) {};
\node[vie, dot] at (2.000,2.170) {};
\node[vie, dot] at (2.880,2.120) {};
\node[vie, dot] at (3.880,2.120) {};
\node[vreg, dot] at (0.000,3.000) {};
\node[vie, dot] at (0.880,2.880) {};
\node[vreg, dot] at (2.000,3.000) {};
\node[vie, dot] at (2.880,2.880) {};
\node[vie, dot] at (3.848,2.924) {};
\node[vie, dot] at (1.120,2.880) {};
\node[vie, dot] at (1.120,2.120) {};
\node[vie, dot] at (1.880,0.120) {};
\node[vie, dot] at (2.830,1.000) {};
\node[vie, dot] at (3.120,2.880) {};
\node[vie, dot] at (3.076,1.848) {};
\node[vie, dot] at (1.880,1.880) {};
\node[vie, dot] at (2.880,1.880) {};
\node[vie, dot] at (2.120,1.880) {};
\node[vie, dot] at (2.170,1.000) {};
\node[vie, dot] at (1.120,1.880) {};
\node[vie, dot] at (0.120,2.120) {};
\node[vie, dot] at (0.880,1.880) {};
\node[vie, dot] at (3.152,2.076) {};
\node[vie, dot] at (3.880,1.880) {};
\node[vie, dot] at (3.152,1.924) {};
\node[vie, dot] at (3.848,1.076) {};
\node[vie, dot] at (3.076,2.152) {};
\node[vie, dot] at (3.924,2.848) {};

\end{tikzpicture}

  \caption{Saída produzida pelo código: 12 elementos de interface (laranja), 19 vértices novos e três buracos
  (hachurados).}
\end{subfigure}
\caption{O exemplo da apresentação de projeto~\cite{slides}, executado pelo código. Na saída, cada cópia de um vértice
foi deslocada um pouco na direção dos elementos que a usam, para que vértices coincidentes fiquem visíveis (``vista
explodida''). A malha é a do arquivo \code{malha2Hyb.txt}.}
\label{fig:exemplo-intro}
\end{figure}

\subsection{Proposta e contribuições deste relatório}

O algoritmo estudado aqui~\cite{slides} parte de uma observação simples, que dá nome a um dos slides de projeto: embora
o problema seja \emph{abrir arestas e criar elementos}, as decisões difíceis acontecem nos \emph{vértices}. Por isso,
a inserção de cada aresta é decomposta em duas aplicações de operadores locais, uma em cada extremidade, e o operador
é escolhido olhando apenas para o estado daquele vértice. Há quatro operadores (\opname{Canônico},
\opname{Split Element}, \opname{Expand Edge} e \opname{Open Hole}), e a ordem das arestas é livre.

Este relatório tem quatro objetivos:
\begin{itemize}
\item explicar o problema e a estrutura de dados IBHM de forma acessível, inclusive o que acontece com seus vetores a cada inserção (\cref{sec:intro,sec:estruturas});
\item descrever a arquitetura do algoritmo, cada operador e o algoritmo de travessia exatamente como estão
implementados na classe \code{IBHMInterfaceElementBuilder}, a base a ser evoluída (\cref{sec:arquitetura,sec:operadores,sec:travessia});
\item verificar experimentalmente o comportamento do código (\cref{sec:verificacao});
\item confrontar a documentação existente (apresentação, comentários do código e artigo da IBHM) com a implementação,
apontando imprecisões e defeitos (\cref{sec:documentacao}).
\end{itemize}
A \cref{sec:conclusao} reúne próximos passos, inclusive observações para a extensão 3D.
